\chapter{User Manual and Deployment Guide} \label{ap:Appendix1}

\section{User Manual}
This section provides step-by-step user operational guides for the three primary roles supported by the \textbf{RecruitSense} platform: Job Seekers, Corporate Recruiters, and Platform Administrators.

\subsection{A.1 User Registration and Authentication}
\subsubsection{A.1.1 Creating a Candidate Account}
\begin{enumerate}
    \item Navigate to the RecruitSense portal at \texttt{http://localhost:5173}.
    \item Click on the \textbf{Sign Up} button in the top navigation bar.
    \item Select the \textbf{Job Seeker} tab.
    \item Enter your Full Name, Email Address, and Password (minimum 8 characters).
    \item Click \textbf{Create Account}. Upon successful registration, you will be automatically logged in and redirected to your Candidate Dashboard.
\end{enumerate}

\subsubsection{A.1.2 Creating an Employer / Recruiter Account}
\begin{enumerate}
    \item Click on \textbf{Sign Up} and select the \textbf{Employer / Recruiter} tab.
    \item Enter your Corporate Email Address (e.g., \texttt{hr@company.com}). The system will automatically check domain validity.
    \item Enter Company Name, Industry, and Office Location.
    \item Submit the registration form. Recruiter accounts with verified corporate domains gain instant job posting privileges.
\end{enumerate}

\subsubsection{A.1.3 Logging In}
\begin{enumerate}
    \item Click \textbf{Log In} on the top navigation bar.
    \item Enter your registered Email and Password.
    \item Click \textbf{Sign In}. The system authenticates credentials, issues a secure 64-character Bearer token, and routes you to your role-specific dashboard.
\end{enumerate}

\subsection{A.2 Job Seeker Operations}
\subsubsection{A.2.1 Profile Setup and Resume Management}
\begin{enumerate}
    \item Navigate to \textbf{My Profile} from the user dropdown.
    \item Enter your professional summary, primary skillset, educational background, and GitHub/LinkedIn portfolio URLs.
    \item Upload your default PDF resume (maximum 5MB).
    \item Click \textbf{Save Changes}.
\end{enumerate}

\subsubsection{A.2.2 Job Discovery and Live Filtering}
\begin{enumerate}
    \item Click on the \textbf{Explore Jobs} tab in the navigation bar.
    \item Use the search bar to enter desired job titles or technology keywords (e.g., ``Full Stack Developer'', ``Python'').
    \item Refine results using sidebar filters: \textit{Job Type} (Full-time, Part-time, Remote, Contract), \textit{Experience Level}, and \textit{Minimum Salary}.
    \item Click on any job card to inspect the comprehensive job description, company information, and required competencies.
\end{enumerate}

\subsubsection{A.2.3 Applying for a Job and Interpreting the ATS Score Modal}
\begin{enumerate}
    \item On the Job Detail page, click the \textbf{Apply Now} button.
    \item Drag and drop your PDF resume into the upload zone or select your pre-uploaded profile resume.
    \item Click \textbf{Submit Application}.
    \item Within seconds, the \textbf{Instant ATS Match Score Modal} will appear:
    \begin{itemize}
        \item \textbf{Overall Match Percentage:} Displays your composite compatibility score (0--100\%) on a circular color-coded gauge.
        \item \textbf{Matched Skills:} Rendered as green badges with checkmarks, indicating competencies successfully extracted from your resume that match the job criteria.
        \item \textbf{Missing Skills (Skill Gap):} Rendered as red warning badges, highlighting required competencies missing from your resume to guide your career upskilling.
    \end{itemize}
\end{enumerate}

\subsubsection{A.2.4 Tracking Application Progress}
\begin{enumerate}
    \item Navigate to \textbf{My Applications} from the top menu.
    \item Review all submitted applications along with their active pipeline status badge (\textit{Applied}, \textit{Screening}, \textit{Shortlisted}, \textit{Interview Scheduled}, \textit{Offer Extended}, \textit{Rejected}).
    \item For applications with status \textit{Interview Scheduled}, click \textbf{View Details} to inspect the scheduled date, time, and video meeting link.
\end{enumerate}

\subsection{A.3 Corporate Recruiter Operations}
\subsubsection{A.3.1 Creating and Publishing a Job Vacancy}
\begin{enumerate}
    \item From the Recruiter Portal, click \textbf{+ Post a Job}.
    \item Enter Job Title, Department, Workplace Type (On-site, Remote, Hybrid), and Location.
    \item Provide a detailed Job Description and outline primary responsibilities.
    \item In the \textbf{Required Skills} field, enter essential technology tags separated by commas (e.g., \texttt{React, Laravel, MySQL, Docker}).
    \item Set compensation parameters and click \textbf{Publish Vacancy}. The job immediately becomes searchable to candidates across the platform.
\end{enumerate}

\subsubsection{A.3.2 Inspecting Ranked Candidate Shortlists}
\begin{enumerate}
    \item Navigate to \textbf{Manage Jobs} and select an active vacancy.
    \item Click \textbf{View Applicants}. The candidate table is automatically rendered in descending order of AI-computed ATS Match Score.
    \item Adjust the \textbf{Minimum ATS Score Filter} slider (e.g., $\ge 80\%$) to automatically filter out under-qualified applicants.
    \item Click on any applicant row to expand the \textbf{Candidate Drawer}, showing extracted resume text, verified matched competencies, and missing skill gaps.
\end{enumerate}

\subsubsection{A.3.3 Managing Pipeline Stages and Scheduling Interviews}
\begin{enumerate}
    \item In the Candidate Drawer or Pipeline Kanban board, select the status dropdown to advance the candidate (e.g., from \textit{Screening} to \textit{Shortlisted}).
    \item When selecting \textit{Schedule Interview}, the \textbf{Interview Drawer} opens:
    \begin{itemize}
        \item Select the Interview Date and Time.
        \item Provide the Video Conference URL (e.g., Google Meet / Microsoft Teams).
        \item Enter special instructions for the candidate.
    \end{itemize}
    \item Click \textbf{Confirm Schedule}. The system automatically updates the application record and dispatches an automated interview invitation to the candidate.
\end{enumerate}

\subsection{A.4 Platform Administrator Operations}
\subsubsection{A.4.1 Corporate Domain Verification Queue}
\begin{enumerate}
    \item Log in to the Admin Governance Console at \texttt{/admin}.
    \item Navigate to \textbf{Employer Verification Queue}.
    \item Review pending company registrations, inspect company website domains and corporate email records.
    \item Click \textbf{Approve} to grant verified employer status or \textbf{Reject} with feedback notes.
\end{enumerate}

\subsubsection{A.4.2 Job Moderation and User Account Management}
\begin{enumerate}
    \item Navigate to \textbf{Job Moderation} to review active or flagged job postings.
    \item Flag, unflag, or permanently delete fraudulent job postings.
    \item Under \textbf{User Management}, search users by email, review activity logs, or suspend abusive accounts with immediate session revocation.
\end{enumerate}

\section{System Deployment and Environment Setup}
To deploy and execute the \textbf{RecruitSense} web platform and AI microservice locally or in a cloud hosting environment, follow the step-by-step instructions below.

\subsection{System Prerequisites}
\begin{itemize}
    \item \textbf{PHP Runtime:} PHP 8.2 or higher with extensions: \texttt{bcmath}, \texttt{curl}, \texttt{mbstring}, \texttt{openssl}, \texttt{pdo\_mysql}, \texttt{xml}.
    \item \textbf{Dependency Manager (PHP):} Composer 2.7+.
    \item \textbf{JavaScript Runtime \& Package Manager:} Node.js 18.x or 20.x LTS and NPM 9.x+.
    \item \textbf{Python Runtime:} Python 3.10 or 3.11 with \texttt{pip} and \texttt{virtualenv}.
    \item \textbf{Database Server:} MySQL Community Server 8.0+.
\end{itemize}

\subsection{Step-by-Step Installation}
\subsubsection{1. Backend Setup (Laravel REST API)}
\begin{verbatim}
cd backend-laravel
composer install
cp .env.example .env
php artisan key:generate
# Configure DB_DATABASE, DB_USERNAME, DB_PASSWORD in .env
php artisan migrate --seed
php artisan storage:link
php artisan serve --port=8000
\end{verbatim}

\subsubsection{2. AI Microservice Setup (Python FastAPI)}
\begin{verbatim}
cd ai-service-fastapi
python -m venv venv
# On Windows:
venv\Scripts\activate
# On Linux / macOS:
source venv/bin/activate
pip install -r requirements.txt
uvicorn main:app --reload --port 5000
# Listening on http://127.0.0.1:5000 (Swagger docs at /docs)
\end{verbatim}

\subsubsection{3. Frontend Setup (React 18 / Vite SPA)}
\begin{verbatim}
cd frontend-react
npm install
# Configure VITE_API_BASE_URL=http://localhost:8000/api in .env
npm run dev
# Accessible at http://localhost:5173
\end{verbatim}

\section{Sample API Contracts}
\noindent \textbf{Analyze Resume Endpoint:} \texttt{POST http://127.0.0.1:5000/analyze-resume}
\begin{verbatim}
Content-Type: multipart/form-data
Form Data:
  resume: [Binary PDF Stream]
  required_skills: "React, Laravel, MySQL, REST API, Tailwind CSS"
  job_description: "Full-Stack Engineer with React and Laravel experience..."

Response (HTTP 200 OK):
{
  "message": "Resume analyzed successfully",
  "similarity_score": 88.45,
  "skill_gap_score": 80.00,
  "matched_skills": ["react", "laravel", "mysql", "rest api"],
  "missing_skills": ["tailwind css"],
  "bonus_skills": ["python", "docker", "git"],
  "total_required_skills": 5,
  "matched_count": 4,
  "character_count": 3420
}
\end{verbatim}